\subsection{Intro to Cryptography}


\content{
        In this chapter, we will be assisting Alice and Bob communicate with each
        other. They are very secretive and do not like their conversations to be
        known by others. 

        Then there's Eve...
}{% presentation
\vspace{2in}
}{% nopresentation
\vspace{2in}
}{% comments
}

\content{
        Let's talk about how we might do this for a second. Suppose that you and
        the person next to you live in different cities (for the purpose of this
        exercise, we all live in a time where messages must be written down and
        taken by a courier to wherever they need to go). You are visiting your
        friend and you decide you want to be able to send secret messages to
        each other, but you don't trust your couriers not to read your message.
        How might you decide to make your messages secret?
}{% presentation
\vspace{3in}
}{% nopresentation
\vspace{3in}
\newpage
}{% comments
}

\subsection{Shift Ciphers}

\content{
        \begin{example}
        A shift cipher is a simple kind of cipher where the alphabet is shifted
        to obscure the original message. For example, if we choose to use a
        shift cipher, and our secret key is 3, then we encode messages according
        to the following: 
        
        \texttt{ABCDEFGHIJKLMNOPQRSTUVWXYZ}\\

        \vspace{-15pt}\hspace{16pt}\texttt{ABCDEFGHIJKLMNOPQRSTUVWXYZ}
        \end{example}
}{% presentation
\vspace{3in}
}{% nopresentation
\vspace{3in}
}{% comments
}

\content{
        \begin{example}
        Use a shift cipher with key 5 to encrypt the message ``We ride at
        dawn''.

        \texttt{ABCDEFGHIJKLMNOPQRSTUVWXYZ}\\

        \vspace{-15pt}\hspace{26pt}\texttt{ABCDEFGHIJKLMNOPQRSTUVWXYZ}
        \end{example}
}{% presentation
\vspace{3in}
}{% nopresentation
\vspace{3in}
}{% comments
}

\content{
        \begin{example}
        The following message was sent to you by a friend, and you know that
        they used a shift cipher with key 8. What was the original message?
        
        \texttt{LWVBFMIBFBPMFNQAP}
        \end{example}
}{% presentation
\vspace{3in}
}{% nopresentation
\vspace{3in}
}{% comments
}


\content{
        How secure is this method? What should we consider a `secure' method?
}{% presentation
}{% nopresentation
}{% comments
knowing what encryption scheme is being used should not make it easier to
decrypt the message
}
